\section{Failure histories and removal of pending sites}
\label{sec:history}

We now use the conditional trial bounds to show that all pending sites
disappear with high probability before the last half of the scales. A
final sweep will then replace every remaining partial potential by its
true value. Throughout this section, the dimension $d\ge3$, the energy
$E\in I_d$, and the coupling $\lambda$ are fixed. Probabilities initially
include the auxiliary randomness used in the construction, and every
cutoff below may depend on $d$.

At the start of a reached scale $n$, a coupled active site has used
exactly its first $n-n_0$ true digits and no digit in its remaining suffix
has been queried. The ordinary rules in Section~\ref{sec:local-sweeps}
and the clearing rules in Section~\ref{sec:assembly} prove this invariant;
Lemma~\ref{lem:conditional-freshness} supplies the conditional law of
the unqueried coordinates. Pending is a conservative status at which the
construction makes no such claim; a pending site can leave that status
only by becoming a committed terminal. In particular, the same invariant
survives a numerical rollback: the affected active sites are marked
pending, while the revealed answers stay in the retained history.

At each scheduled trial start reached after a finite history produced by
these rules, including histories with failed trials, these checks verify
the random-input hypotheses of Propositions~\ref{prop:ordinary-sweep}
and~\ref{prop:clearing-trial}. The ordinary argument also proves, by
successive conditioning, that auxiliary retries terminate almost surely.
Each clearing trial has a finite heat limit. Its ridge choices and the
diagonal reset after a stage are measurable functions of the known
numerical state, so neither exposes further randomness. At fixed $m$ there are finitely
many scales, tiles, and snapshot components, so the construction reaches each
subsequent scheduled trial after an almost surely finite number of
operations. We use these conditional conclusions below; the numerical
bounds along committed states are the pathwise conclusion of
Proposition~\ref{prop:deterministic-induction}.

\subsection{Birth labels and chronological probability bounds}

A \emph{birth} is a first tile-trial failure in a one-digit sweep or a
failed clearing trial. Give a one-digit sweep from scale $j-1$ to $j$ the
post-sweep index $j$. The late-failure bound below treats the final residual
sweep separately. A birth
label consists of this index, its type, and its tile or snapshot-component
anchor.  Tile and clearing labels are distinct even when their physical
anchors coincide. Each birth assigns its label to every site covered by
that failure's marking rule, replacing any label the site already carried.
At the end of each clearing stage, including the last stage before the
cutoff check, every pending site thus carries its latest birth label.

There is a deterministic ordering of all possible labels: increasing scale,
the prescribed tile order within each sweep, and then increasing site
order for possible clearing anchors.  At an anchor absent from the frozen
snapshot, or belonging to a component too large to be attempted, insert a
skipped slot.  Auxiliary retries create no further birth labels.
Use one designated site as the label anchor for every one-digit tile.
After a global abort, continue this bookkeeping list by recording all
remaining slots as skipped. The list of
possible labels is deterministic even though eligibility of a particular
clearing anchor and the time spent on earlier retries depend on history.

Let $B_a$ be the indicator of a birth at a potential label $a$, with value
zero at a skipped slot, and let $\mathcal F_a^-$ be the retained history
just before that slot. In the operation filtration of
Lemma~\ref{lem:conditional-freshness}, the time at which this slot is
reached is a stopping time: the preceding queries, trial outcomes, and
decisions to schedule or skip a slot are all recorded. It is almost surely
finite by the retry bound above. Thus $\mathcal F_a^-$ is the corresponding
stopped history, and whether the slot is scheduled is
$\mathcal F_a^-$-measurable. The uniform conditional trial estimates imply,
for some fixed $c>0$ and sufficiently large $n_0$,
\begin{equation}
 \EE(B_a\mid\mathcal F_a^-)
 \le p_{j(a)},\qquad p_j=\exp(-4^{cj}).
 \label{eq:history-birth}
\end{equation}
Here the exponent has been decreased to absorb fixed constants and the
one-unit difference between sweep and post-sweep indices.  Indeed,
writing $a_\lambda=(\lambda^2/3)4^{n_0}$ gives
\begin{equation}
 \tfrac14<a_\lambda\le1,\qquad
 t_j=a_\lambda(4^{-j}+4^{-m}),\qquad
 \tfrac14 4^{-j}<t_j\le2\,4^{-j}\quad(j\le m).
 \label{eq:history-scale-comparison}
\end{equation}

For any distinct, deterministically specified labels $a_1<\cdots<a_s$,
\begin{equation}
 \PP(B_{a_1}=\cdots=B_{a_s}=1)
 \le\prod_{b=1}^s p_{j(a_b)}.
 \label{eq:history-joint}
\end{equation}
To prove this, the product of the first $s-1$ indicators is measurable
before the last slot. Apply~\eqref{eq:history-birth} to the last indicator
and repeat backwards. The stopped histories include every query made
during the almost surely finite intervening retries. Hence their random
durations do not affect the chronological argument.
Equation~\eqref{eq:history-joint} does not assert independence of births.
Adaptive anchors will be handled by a union over deterministic possible
labels.

\subsection{Witnesses for a surviving pending site}

The remaining task is to exploit the small conditional probabilities of
births.  Pending sites can survive a clearing stage without a new failure
only by belonging to a component too large to clear.  Such a component
provides many separated pending sites at a much earlier scale.

An eligible component at scale $j$ has at most $t_j^{-\nu}$ sites.  Its
$R_j$-connectedness and the inspection radius imply that every site marked
by its failed trial lies at distance at most
\begin{equation}
 C t_j^{-\nu}R_j+\ell_j\le C4^{(8+\nu)j}
 \label{eq:history-mark-radius}
\end{equation}
from its anchor.  The same bound holds for a tile birth, whose marks lie in
a tile of side $4^{6(j-1)}$.

After clearing stage $j$, every pending site is either a mark of a failure
at that stage or belongs to a huge snapshot component, meaning one with
more than $t_j^{-\nu}$ vertices.  This includes pending marks made during
the preceding sweep, since they were present when the snapshot was taken.
Eligible components disappear on success and receive a new label on
failure. Separation and the pivot-locality argument in
Section~\ref{sec:assembly} prevent another component's trial from touching
a huge component.

We also use the following persistence observation.  If a site is pending
at the end of stage $j$ and its latest birth index is at most $r<j$, then it
was pending at the end of stage $r$.  It could not have become coupled in
between, and if it had become terminal it could never have been marked
again.

\begin{proposition}[Pending-site bound]
\label{prop:witness-bound}
There is a constant $c_{\rm h}>0$ such that, after increasing the fixed lower
cutoff $n_0$ if necessary, for every $n_0\le j\le\lfloor m/2\rfloor$ and
every torus site $x$,
\[
 \PP(x\text{ is pending at the end of stage }j)
 \le\exp(-4^{c_{\rm h}j}).
\]
All constants are uniform in $m,E,\lambda$ in the stated ranges.
\end{proposition}

\begin{proof}
Put
\[
 \alpha_{\rm h}=\nu/(100d),\qquad \beta_{\rm h}=0.2\nu,\qquad
 W_j=4^{10j},\quad r=\lfloor\alpha_{\rm h} j\rfloor,\quad
 k_j=\lfloor4^{\beta_{\rm h} j}\rfloor.
\]
The radius $W_j$ will contain both the possible reach of a stage-$j$
failure and the connected exploration of a huge component near the
root. Passing to the much earlier scale $r$ makes a ball of radius
$3W_r$ small compared with that component. We will retain only $k_j$
of the separated older sites obtained this way. The estimates below
verify these three uses of the choices.

We define a deterministic family of finite rooted trees, called witnesses.
A leaf of a scale-$j$ witness rooted at $x$ specifies one birth label of
index $s$ with $n_0<s\le j$ and $s>r$, whose anchor is within distance
$W_j/2$ of $x$.
Alternatively, when $r\ge n_0$, an internal vertex specifies $k_j$ child
centers within distance $W_j/2$ of $x$, separated from each other by more
than $3W_r$, and a scale-$r$ witness at each child.  The family at scale
$n_0$ is empty, since there are initially no pending sites.  The scales
strictly decrease down a branch, so every tree is finite.

\paragraph{Extracting a witness and separating its leaf labels.}
We first show that a pending site yields a witness all of whose leaf births
occurred.  A site newly marked by a clearing failure has a stage-$j$ birth
nearby, giving a leaf by~\eqref{eq:history-mark-radius}.  Otherwise it lies
in a huge snapshot component.  Explore a connected set of
\[
 M_j=\lfloor t_j^{-\nu}\rfloor+1\ge c_1 4^{\nu j}
\]
distinct vertices starting at $x$.  Each is within distance
$M_jR_j\le C4^{(8+\nu)j}$ of $x$.  If one has latest birth index $s>r$,
its birth anchor is within
$C4^{(8+\nu)j}+C4^{(8+\nu)s}<W_j/2$ of $x$, once $n_0$ is large.
This again gives a leaf.

If no such latest birth exists, all explored sites were pending at end
$r$, by persistence.  Necessarily $r\ge n_0$.  A ball of radius $3W_r$
contains at most
\[
 C W_r^d=C4^{10dr}\le C4^{0.1\nu j}
\]
sites.  Greedy packing of the explored set therefore supplies at least
$c_2 4^{0.9\nu j}$ centers with the required separation, more than $k_j$
for sufficiently large $n_0$.  Choose $k_j$ and recursively extract their
witnesses.  Their centers lie within $W_j/2$ of $x$.

Every birth anchor in a scale-$j$ witness is within $W_j$ of its root.
For a leaf this is immediate.  For an extracted branch it follows from
\[
 C4^{(8+\nu)j}+W_r<W_j.
\]
For the entire deterministic witness family the same conclusion follows
from $W_j/2+W_r<W_j$.  Thus different child witnesses have disjoint sets
of possible physical birth anchors: their support balls have radius $W_r$
and centers more than $3W_r$ apart.  Induction shows that all leaf birth
labels in a witness are distinct.  This is the reason for spatial
separation; no spatial independence is needed.

\paragraph{Summing deterministic witness weights.}
First fix a deterministic possible witness tree, including all its
centers and leaf labels. Its leaves have distinct labels, so
\eqref{eq:history-joint} bounds the probability that all its leaves occur
by the product of their $p_s$'s. Assign this product as the tree's weight.
Let $Z_j$ be the supremum over roots of the sum of these weights over
all allowed scale-$j$ witnesses, and put $Z_{n_0}=0$. The deterministic
family is finite at every scale. A union bound now bounds the pending
probability by $Z_j$. The adaptive extraction above proves only existence
of a member of this family; we do not condition on the tree extracted
from a realized component. A related expansion into deterministic trees
of separated lower-scale events appears in Sznitman's treatment of
cascading events~\cite[Definition~3.1 and proof of Theorem~3.4]{Sznitman2012}.
Here the chronological conditional bound~\eqref{eq:history-joint}
supplies the product weight for distinct birth labels.

There are at most $CjW_j^d$ possible leaf labels, since
there are at most $j$ scales, two types, and $CW_j^d$ anchor positions.
In counting branches, order the child centers and discard their separation
constraint to obtain an upper bound.  Hence
\begin{equation}
 Z_j\le CjW_j^d\exp\{-4^{c\max(n_0,r)}\}
       +(CW_j^d Z_r)^{k_j},
 \qquad r=\lfloor\alpha_{\rm h} j\rfloor.
 \label{eq:history-recursion}
\end{equation}
The second term is absent when $r<n_0$, and vanishes at $r=n_0$.

\paragraph{Closing the numerical recursion.}
Choose
\begin{equation}
 c_{\rm h}=\min\{c\alpha_{\rm h}/4,\nu/20\}>0.
 \label{eq:history-exponent}
\end{equation}
Then $c_{\rm h}<c\alpha_{\rm h}$ and
$\beta_{\rm h}+c_{\rm h}\alpha_{\rm h}-c_{\rm h}\ge.15\nu>0$.
We prove $Z_j\le e^{-4^{c_{\rm h}j}}$ by induction. The leaf term in
\eqref{eq:history-recursion} requires care when $\alpha_{\rm h}$ is tiny.
If $r<n_0$, then
\[
 j<(n_0+1)/\alpha_{\rm h},\qquad
 c_{\rm h}j\le c(n_0+1)/4.
\]
Its logarithm is at most
\[
 \log(Cj)+10dj\log4-4^{cn_0}
 \le -4^{c_{\rm h}j}-\log2
\]
after one fixed increase of $n_0$: the entropy is
$O_d(n_0/\alpha_{\rm h})$, whereas $4^{cn_0}$ dominates both that
quantity and $4^{c(n_0+1)/4}$. This treats the entire, possibly very
long, nonrecursive range. If $r\ge n_0$, use
$j<(r+1)/\alpha_{\rm h}$ and
$c_{\rm h}j\le c(r+1)/4$ instead. Now $4^{cr}$ dominates the
entropy $O_d((r+1)/\alpha_{\rm h})$ and the target exponent
$4^{c(r+1)/4}$, uniformly for $r\ge n_0$. The same half-bound for
the leaf term follows. These thresholds depend on $d$ but not on $m$.

For the second term, assume $r\ge n_0$ and $Z_r>0$; otherwise it
vanishes.  A single fixed increase of $n_0$
ensures, for every integer $u\ge n_0$,
\begin{equation}
 4^{c_{\rm h}u}\ge
 2\left(\log C+10d\log4\,\frac{u+1}{\alpha_{\rm h}}\right).
 \label{eq:history-initial-entropy}
\end{equation}
Such a choice is possible because exponential growth dominates the fixed
linear function on the right.  Since $j<(r+1)/\alpha_{\rm h}$, the induction
hypothesis now gives
\[
 \log(CW_j^d Z_r)\le-\tfrac12 4^{c_{\rm h}r}.
\]
Using $k_j\ge4^{\beta_{\rm h} j}/2$ and $r\ge\alpha_{\rm h} j-1$, we obtain
\[
 (CW_j^d Z_r)^{k_j}
 \le\exp\{-\tfrac14 4^{-c_{\rm h}}4^{(\beta_{\rm h}+c_{\rm h}\alpha_{\rm h})j}\}
 \le\tfrac12 e^{-4^{c_{\rm h}j}}.
\]
Together with the first half-bound this closes the induction.
Condition~\eqref{eq:history-initial-entropy} explicitly includes the first
recursive scales, when the parent index can be nearly $n_0/\alpha_{\rm h}$;
it requires no dependence of $n_0$ on $m$.

Finally, these geometric counts are valid on the torus.  Since
$j\le\lfloor m/2\rfloor$, every support radius is at most
$W_j\le4^{5m}$, whereas the side is $4^{20m}$.  All exploration and
separation estimates therefore use the same uniform $d$-dimensional
ball bounds as on the lattice.
\end{proof}

\subsection{The cutoff and the last true-tail sweep}

We can now finish the finite-volume construction.  Write
$J=\lfloor m/2\rfloor$.  Proposition~\ref{prop:witness-bound} and the
volume $N=4^{20dm}$ give
\begin{equation}
 \PP(\text{some site is pending at end }J)
 \le\exp\{20dm\log4-4^{c_{\rm h}J}\}\longrightarrow0.
 \label{eq:history-cutoff}
\end{equation}
Abort if this event occurs.  Otherwise every active site is coupled.
Thereafter no clearing is scheduled, so a new pending site would have no
later route to a committed terminal before the final sweep. We therefore perform
ordinary sweeps without clearing and abort at the first first-trial failure.
There are at most $(m+1)N$ possible tile trials in
these stages, including the final sweep.  Give a trial after a previous
abort failure indicator zero.  Successive conditioning and the uniform
trial bound yield
\begin{equation}
 \PP(\text{a later ordinary failure before abort})
 \le C(m+1)N\exp(-4^{cJ})\longrightarrow0,
 \label{eq:history-late-failure}
\end{equation}
decreasing $c$ once more if necessary.  No independence between the
late trials is asserted.

These two errors are also small uniformly over every $m\ge4n_0$,
once $n_0$ is sufficiently large. Indeed, with
$q=\min(c,c_{\rm h})>0$ and $a=20d\log4$, their sum is at most
\begin{equation}
 [1+C(m+1)]\exp\{am-4^{q(m/2-1)}\}.
 \label{eq:history-uniform-abort}
\end{equation}
Choose a dimension-dependent $M$ so large that for all $m\ge M$,
$4^{q(m/2-1)}\ge(a+2)m$ and $1+C(m+1)\le e^m$.
Then \eqref{eq:history-uniform-abort} is at most $e^{-m}$ whenever
$m\ge4n_0\ge M$. The scale comparison
\eqref{eq:history-scale-comparison} makes this estimate uniform in
$a_\lambda\in(1/4,1]$ and $E\in I_d$ as well.

Consider a retained history that reaches the start of scale $m$ without
abort. Every active site is still coupled, so its true residual $\xi_i$
has not been queried. Lemma~\ref{lem:conditional-freshness} gives its
uniform law with variance $s_m=\eta$, conditionally independent between
sites. Reaching this stage is decided by retained information; the law
does not condition on eventual success of the final sweep. For that sweep
take
\[
 t=2\eta,\qquad w=\theta=\tfrac12,\qquad
 \xi_{0,i}=\xi_i/\sqrt w=\sqrt2\,\xi_i.
\]
Then, conditional on its pre-query history,
\[
 \EE \xi_{0,i}=0,\qquad \EE \xi_{0,i}^2=t,\qquad |\xi_{0,i}|\le\sqrt{3t}.
\]
These are the bounded-input moment hypotheses of
Proposition~\ref{prop:ordinary-sweep}.  The entire residual, rather than
one digit, may therefore be consumed in this final sweep.  If
$k=w^{1/4}$ and $r_{\rm bit}=k^2r_{\rm phys}$, the physical diagonal
update is exactly
\[
 k^2(\xi_{0,i}+t r_{\rm bit})
 =\sqrt w\,(\xi_i/\sqrt w+t\sqrt w\,r_{\rm phys})
 =\xi_i+\eta r_{\rm phys}.
\]
It inserts the full true residual and leaves the artificial term
$-\eta r_{\rm phys}$ from the previous $-2\eta r_{\rm phys}$.
The reserve normalization is $\sqrt{\theta t}=\sqrt\eta$.
No reset is needed after this sweep.

\begin{proof}[Completion of the proof of Proposition~\ref{prop:certificate}]
Initialization and Proposition~\ref{prop:deterministic-induction} provide
the committed-state invariants throughout the construction.
Lemma~\ref{lem:conditional-freshness}, together with the ordinary and
clearing query rules verified in Sections~\ref{sec:local-sweeps}
and~\ref{sec:assembly}, supplies all conditional random-input hypotheses,
including after failed trials. Preliminary auxiliary retries
terminate almost surely.  By~\eqref{eq:history-cutoff}
and~\eqref{eq:history-late-failure}, with probability $1-o_m(1)$ the
cutoff finds no pending site and every subsequent sweep, including the
whole-tail sweep, succeeds.

On that event all terminal sites have their exact true potentials, and the
last sweep makes every remaining active potential true as well.  The
active set $A$ has cardinality at least $N/2$ by the deterministic
induction.  Its final frozen coefficients satisfy
$|r_i|\le C$ and $c\le\Im r_i\le C$; the endpoint diagonal-lag bound
of the last ordinary sweep gives an active diagonal imaginary part at
least a fixed positive constant.  The final active inverse is precisely
the Schur inverse of
\[
 H_m-E-\eta\diag_A r.
\]
The true terminal principal block is invertible on the probability-one
event established in Section~\ref{sec:state}, and the active Schur block has strictly
negative definite imaginary part before inversion.  Thus the full inverse
exists and has the asserted active diagonal entries.

The analytic interpolation exponents remain $4/3,2/3$ in all dimensions:
Section~\ref{sec:initialization} proves a deliberately weaker Fourier
bound requiring only $d\ge3$. The common choice of $\nu$ leaves
$d\nu=10^{-5}\delta$, below every analytic gain consumed by packing.
The additional factor $d$ in the history scale ratio is needed for
the count of spatial witnesses in Section~\ref{sec:history}.

All numerical lower bounds on $n_0$, including
\eqref{eq:history-initial-entropy}, are fixed independently of $m$,
$a_\lambda\in(1/4,1]$, and $E\in I_d$. For every starting index above
this dimension-dependent cutoff, the certificate probability tends to
one as $m\to\infty$ and is at least $1/2$ uniformly for all
$m\ge4n_0$ after a further fixed increase of the cutoff. This proves
the certificate with its stated uniform constants.
\end{proof}

Although $A$ and $r$ may depend on both true and auxiliary variables, the
true potential vector has its prescribed iid law as the unconditional
marginal of the enlarged space. Each joint success witnesses the assertion
in the true variables that a qualifying certificate exists for $H_m$.
Section~\ref{sec:spectral} proves that this feasibility event is Borel.
Its probability under the true product law is at least the joint success
probability. The spectral averaging is then unconditional under that law.

\subsection{An intrinsic positive disorder threshold}

In this subsection only, fix $d\ge4$. We specify its threshold by
finite matrices and a fixed success probability. This avoids leaving
an arbitrary choice among the many smallness constants in the
preceding estimates. For $d=3$ we retain the existential threshold
in Proposition~\ref{prop:certificate}.

\begin{definition}[Finite-box threshold]
\label{def:threshold}
For the fixed dimension $d\ge4$, let $n_*(d)$ be the least integer
$k\ge10$ with the following property. For every pair of integers
$n_0\ge k$ and $m\ge4n_0$, every $u\in(1/4,1]$, and every
$E\in I_d=-2d+(1/200,1/100)$, set
\[
 L=4^{20m},\qquad N=L^d,\qquad
 \lambda(n_0,u)=\sqrt{3u}\,2^{-n_0},\qquad \eta=u4^{-m}.
\]
Let $H_{m,n_0,u}$ be the true periodic negative-hopping Anderson matrix
on $(\ZZ/L\ZZ)^d$ with independent uniform $[-1,1]$ variables and
coupling $\lambda(n_0,u)$. With probability at least $1/2$, there must
exist a set $A$ and coefficients $r=(r_i)_{i\in A}$ such that
\begin{equation}
 \begin{gathered}
 |A|\ge N/2,\qquad |r_i|\le k,\qquad \Im r_i\ge1/k\quad(i\in A),\\
 G_A^{\rm cert}:=
 (H_{m,n_0,u}-E-\eta\diag_A r)^{-1}\ \text{exists},\qquad
 \Im(G_A^{\rm cert})_{ii}\ge1/k\quad(i\in A).
 \end{gathered}
 \label{eq:threshold-criterion}
\end{equation}
The probability here is taken only over the true site variables.
Define
\begin{equation}
 \lambda_d=\sqrt3\,2^{-n_*(d)}.
 \label{eq:threshold-lambda}
\end{equation}
\end{definition}

\begin{proposition}[Existence and measurability of the threshold]
\label{prop:threshold}
The feasibility event in \eqref{eq:threshold-criterion} is Borel, and
the set of integers in Definition~\ref{def:threshold} is nonempty.
Consequently $\lambda_d>0$. For every fixed $0<\lambda<\lambda_d$,
every $m\ge4n_0$, and every $E\in I_d$, the true operator at coupling
$\lambda$ satisfies \eqref{eq:threshold-criterion} with
$k=n_*(d)$ and certificate probability at least $1/2$.
\end{proposition}

\begin{proof}
First fix $k,n_0,m,u,E$ and $A$. Write
\[
 B_A(\omega,r)=H_{m,n_0,u}(\omega)-E-\eta\diag_A r.
\]
The product of the potential cube $[-1,1]^N$ and
\[
 \mathcal R_A(k)=\{r\in\CC^A:|r_i|\le k,\ \Im r_i\ge1/k\}
\]
is compact. For each positive integer $q$, impose on this product
the additional conditions
\[
 |\det B_A(\omega,r)|\ge1/q,
 \qquad \Im(B_A(\omega,r)^{-1})_{ii}\ge1/k\quad(i\in A).
\]
The inverse is continuous on the determinant-bounded set, so the
resulting set is compact. Its projection onto the potential coordinates
is compact. Taking the countable union over $q$ and the finite union
over $A$ with $|A|\ge N/2$ gives exactly the feasibility event, hence a
Borel set. This argument uses the full finite matrix, without any
extra condition on a Schur denominator.

For existence, all constants and starting-index requirements in
Proposition~\ref{prop:deterministic-induction} and the witness proof
depend only on $d$ and on the fixed window. They are uniform in
$u\in(1/4,1]$ because $t_n=u(4^{-n}+4^{-m})\asymp4^{-n}$.
Equation~\eqref{eq:history-uniform-abort} makes the total failure
probability at most $1/2$ simultaneously for every $m\ge4n_0$ once
$n_0$ exceeds one fixed cutoff. On success the certificate has uniform
bounds $|r_i|\le C_r$, $\Im r_i\ge c_p$, and
$\Im G_{ii}^{\rm cert}\ge c_g$, with all three constants positive
where appropriate. Choose an integer $k\ge10$ above that cutoff and
above $C_r,c_p^{-1},c_g^{-1}$. Then the constructed success event implies
\eqref{eq:threshold-criterion} for every allowed set of parameters.
Although the construction also uses auxiliary randomness, its success
implies the Borel feasibility event in the true variables alone.
The latter event therefore has probability at least the joint success
probability, which is at least $1/2$. Thus such $k$ exist.

Finally, if $0<\lambda<\sqrt3\,2^{-n_*(d)}$, set
\[
 n_0=\left\lfloor\log_4(3/\lambda^2)\right\rfloor,
 \qquad u=(\lambda^2/3)4^{n_0}.
\]
Then $n_0\ge n_*(d)$, $1/4<u\le1$,
$\sqrt{3u}\,2^{-n_0}=\lambda$, and $s_m=u4^{-m}$.
Definition~\ref{def:threshold} applies directly. Universal quantification
over $u$ and $E$ refers to their individual finite-box probabilities;
it does not require a common probability-one event over those parameters.
\end{proof}

The existence proof supplies no effective numerical upper bound for
$n_*(d)$. The spectral readout uses the fixed probability
$1/2$ and the certificate bounds in \eqref{eq:threshold-criterion} at the
torus and damping scales prescribed in Definition~\ref{def:threshold}.
Its comparison holds for every
qualifying pair $(A,r)$, so no measurable choice of a certificate and
no independence between $A$ and the true variables are required.
